\section{Discussion}

\subsection{LLMs as a tool for literature synthesis: scale, semantic reach, and verification}

The corpus-building process illustrates both the value and the limits of LLM-assisted review. Language models can help screen and organize evidence across more publications than reviewers could process manually at the same speed. Used alongside database searches and explicit eligibility criteria, they can support consistent classification of topics, modeled entities, and methodological roles in titles and abstracts. Their semantic capabilities may also distinguish related themes that exact keyword searches or fixed dictionaries overlook, improving thematic discrimination beyond surface-level matching.

These advantages do not guarantee accuracy. Classification quality depends on corpus coverage, operational definitions, model and prompt stability, and the evidence available in each record. Screening and coding decisions therefore remain open to audit. LLM assistance can shift rather than remove human work: reviewers must examine ambiguous cases, resolve disagreements, verify citations and extracted claims, and document inclusion and exclusion decisions. This verification burden can be substantial at scale. Risk-based review offers one practical response: double-code stratified samples and high-impact or uncertain cases, report error types and agreement, and preserve prompts, model versions, and decision logs.

\subsection{How LLMs can be embedded in ABM research}

LLMs can contribute to ABM at several distinct stages, and these roles should be distinguished. As an internal agent mechanism, an LLM can interpret context, generate candidate actions, support planning or dialogue, and mediate adaptation. This offers a direct route to language-mediated behavior, but makes decisions sensitive to prompts, model updates, sampling, and context construction. As an external research assistant, an LLM can help translate theory into candidate rules or code, generate tests and scenarios, extract behavioral evidence from documents, or provide natural-language interfaces for inspecting a model. These supporting uses may accelerate research without making the LLM the causal decision mechanism of the simulated agents.

A sound research design should specify the LLM’s role and where it enters the modeling workflow. For LLM-driven agents, authors should define the agent population and information access, distinguish generated behavior from hard-coded constraints, and test sensitivity to model, prompt, seed, and temperature. Validation should compare individual choices and interactions with empirical or experimental benchmarks, then assess whether aggregate outcomes and policy responses remain credible. For supporting applications, researchers should document and independently verify the provenance of generated code, labels, and synthetic data. In either case, LLMs should remain components of a theory-led, empirically constrained workflow—not substitutes for behavioral theory, calibration, or validation.

\subsection{Collaboration networks and the development of ABM communities}

The increasingly modular coauthorship structure suggests a dual agenda. Specialized communities provide the depth needed to develop domain-specific theory, data, and validation practices. At the same time, weak ties between communities can slow the circulation of reusable mechanisms and standards. Future research should examine not only who collaborates, but how methods travel: whether cross-community teams, shared benchmarks, open model repositories, and replication projects connect epidemiology, transport, environmental systems, and social dynamics. Longitudinal analyses could test whether LLM-enabled ABM papers bridge existing groups or form a distinct cluster, and whether such positions correspond to methodological diffusion rather than publication growth alone.

These inferences should be treated cautiously. Community assignments depend on author-name disambiguation, database coverage, network thresholds, and the resolution and time window of the detection algorithm; coauthorship is also an imperfect proxy for intellectual exchange. Follow-up work should test robustness across network specifications and combine coauthorship with complementary evidence, including shared citations, software reuse, institutional ties, and cross-domain publications. The goal is not to eliminate disciplinary boundaries, but to build translation infrastructure—interoperable model descriptions, shared behavioral benchmarks, transparent data and code, and recurring replication exercises—so that specialized groups can compare findings without erasing meaningful differences in theory and context.